\section{Experimental details}\label{app:details}

\paragraph{Training-free operators.}
All operators share the score $s=s_1/c_1+\beta s_3/c_3$ (Section~4) with
$\beta\in\{0,0.03,0.1,0.3,1,3,10\}$ or the three-hop term alone; the five-hop
variants add $\gamma s_5/c_5$ with $\gamma\in\{0.03,0.1,0.3,1,3\}$. The grids are:
normalisation exponent $\alpha\in\{0.3,0.5,0.7\}$ (Coat), $\{0.5,0.7\}$
(Yahoo!\,R3) and $\{0.3,0.5\}$ (KuaiRec); clip
$\tau\in\{0.01,0.02,0.05,0.1,0.2\}$ (Coat), $\{0.1,0.2\}$ (Yahoo!\,R3) and
$\{0.05,0.1,0.2\}$ (KuaiRec); imputation ridge $\lambda\in\{1,5,20\}$ (Coat),
$\{20\}$ (Yahoo!\,R3) and $\{5,20\}$ (KuaiRec); additive or low-rank
(rank 32) imputation; control-variate weight $\in\{0,0.25,0.5,1\}$ on Yahoo!\,R3
and $1$ elsewhere; and degree source $\in\{\hat Y,\ \text{cross-fitted }W\}$ for
the debiased operators. The degree source is a searched hyper-parameter; both
options satisfy Remark~\ref{rem:xfit}, and only $\hat Y$-degrees make the
certificates of Theorem~\ref{thm:robust} informative. EASE uses
$\lambda\in\{10,100,1000\}$; GF-CF uses the ideal low-pass weight
$\in\{0,0.3,1\}$ with $64$ or $256$ singular vectors. On the DR graph, EASE and
GF-CF search the same nuisance grid as DR adjacency. The propensity model is
the one shipped with Coat and a logistic exposure model on user and item
exposure rates elsewhere; on Yahoo!\,R3 it was selected on validation data
against an outcome-dependent Naive-Bayes model. All nuisances are cross-fitted
over ten random folds of pairs.

\paragraph{Trained models.}
Embeddings have 64 dimensions (Coat: 32 or 64) and LightGCN-type models two
layers. The grid is learning rate $\{10^{-2},3\cdot10^{-3},10^{-3}\}$ $\times$ weight
decay $\{10^{-5},10^{-4},10^{-3}\}$ on Coat and Yahoo!\,R3 and
$\{10^{-2},3\cdot10^{-3}\}\times\{10^{-5},10^{-4}\}$ on KuaiRec, plus
$\gamma\in\{0.1,0.2\}$ for PDA and $r\in\{0.3,0.7\}$ for r-AdjNorm. BPR models draw
one uniformly sampled negative per positive. Batch sizes are 256 (Coat), 1{,}024
(Yahoo!\,R3) and 32{,}768 (KuaiRec, with 2M sampled pairs per epoch for the
pointwise losses); training runs for at most 100, 60 and 30 epochs with
patience 5, 5 and 3 on the mean validation score. Because the training log is
the same for every split, each configuration is trained once and evaluated
after every epoch on the validation part of every split; each split then
selects its own configuration and epoch. The split-0 selection is re-trained
with three (KuaiRec: two) further seeds to measure seed variance.
Table~\ref{tab:configs} lists the number of configurations each method was
selected from.

\begin{table}[h]\centering\footnotesize
\caption{Number of configurations searched per method (per split).}\label{tab:configs}
\begin{tabular}{lccc}
\toprule
Method & Coat & Yahoo!\,R3 & KuaiRec\\
\midrule
Popularity (non-personalised) & -- & 1 & --\\
Imputation only (non-personalised) & -- & 2 & --\\
MF (trained) & 18 & -- & --\\
IPS-MF (trained) & 18 & -- & --\\
DR-MF (trained) & 18 & -- & --\\
LightGCN (trained) & 18 & -- & --\\
DR-LightGCN (trained) & 18 & -- & --\\
MF (trained) & 18 & -- & --\\
PDA (trained) & 36 & -- & --\\
LightGCN (trained) & 18 & -- & --\\
r-AdjNorm (trained) & 36 & -- & --\\
NAVIP (trained) & 18 & -- & --\\
linear LightGCN / r-AdjNorm (training-free) & -- & 16 & --\\
EASE (training-free) & -- & 3 & --\\
GF-CF (training-free) & -- & 4 & --\\
IPS adjacency (NAVIP-style) (training-free) & -- & 32 & --\\
IPS + walk correction (training-free) & -- & 32 & --\\
EASE on DR graph (training-free) & -- & -- & --\\
GF-CF on DR graph (training-free) & -- & -- & --\\
DR adjacency (training-free) & -- & 256 & --\\
\textbf{DRUP} (training-free) & -- & -- & --\\
DR adjacency, 5 hops (training-free) & -- & -- & --\\
\textbf{DRUP}, 5 hops (training-free) & -- & -- & --\\
\bottomrule
\end{tabular}

\end{table}

\paragraph{Evaluation and tests.}
Coat splits each user's candidates (30\% validation), Yahoo!\,R3 and KuaiRec
split the test users (30\% validation), over 10, 5 and 5 random splits.
Users without a positive candidate are skipped. For a pair of methods, each
test user's metric difference is averaged over the splits in which the user
is in the test part, and we report the mean difference over users with a
95\% $t$-interval and the paired $t$-test, Holm-adjusted over all comparisons
of a dataset. The Wilcoxon signed-rank test gives the same conclusions except for the
five-hop DRUP against DR comparison on KuaiRec ($p=0.66$ against 0.011). The
corrected resampled $t$-test over split means~\cite{nadeau2003} is far more
conservative with five or ten splits and is reported in the result files.

\paragraph{Hardware and code.}
All experiments ran on four CPU cores with 16\,GB of memory, without a GPU.
Code, configurations and result files are available at the repository given
on the title page.
